\documentclass[11pt]{amsart}
\usepackage{amssymb}

\newtheorem{theorem}{Theorem}[section]
\newtheorem{lemma}[theorem]{Lemma}
\newtheorem{proposition}[theorem]{Proposition}
\newtheorem{corollary}[theorem]{Corollary}
\theoremstyle{definition}
\newtheorem{definition}[theorem]{Definition}
\theoremstyle{remark}
\newtheorem*{remark}{Remark}

\DeclareMathOperator{\ord}{ord}
\DeclareMathOperator{\Fix}{Fix}

\title{Involutions in groups of even order: a counting note}
\author{Marta K. Oyelaran}
\address{Department of Mathematics, Larkfield College, Ashbury, UK}
\email{m.oyelaran@larkfield.ac.uk}
\date{}

\begin{document}

\begin{abstract}
We give a short self-contained account of the classical fact that a finite group of even
order contains an element of order $2$. The argument is the usual pairing of an element with
its inverse, organised so that it yields the slightly sharper statement that the number of
elements of order $2$ is odd. We also record two applications.
\end{abstract}

\maketitle

\section{Introduction}

The statement that every finite group of even order contains an element of order $2$ is the
case $p=2$ of Cauchy's theorem, and it is often the first nontrivial fact about finite groups
that students meet. The standard proof pairs each element $g$ with its inverse $g^{-1}$: the
elements that are not paired with themselves come in pairs, so the elements with $g=g^{-1}$
are even in number, and since the identity is one of them there must be another.

The purpose of this note is to write this argument out carefully, separating the
combinatorial input (a parity statement about involutions of a finite set,
Lemma~\ref{lem:parity}) from the group-theoretic input (Lemmas~\ref{lem:div}--\ref{lem:ord2}).
Our main result is Theorem~\ref{thm:main}. In Section~\ref{sec:app} we describe groups in which
every element is its own inverse and recall a theorem of Frobenius.

Throughout, $G$ denotes a finite group with identity $e$, written multiplicatively, and $|G|$
is its order. For an integer $m\ge 0$ we write $g^m$ for the $m$-fold product, with $g^0=e$.

\section{Orders of elements and a parity lemma}

\begin{definition}\label{def:ord}
Let $g\in G$. The \emph{order} of $g$, written $\ord(g)$, is the least positive integer $n$
such that $g^n=e$. Since $G$ is finite, two of the powers $g, g^2, g^3,\dots$ coincide, so
such an $n$ exists.
\end{definition}

\begin{lemma}\label{lem:div}
Let $g\in G$ and let $m\ge 0$ be an integer. Then $g^m=e$ if and only if $m$ divides $\ord(g)$.
\end{lemma}

\begin{proof}
Put $n=\ord(g)$. If $n\mid m$, say $m=kn$, then $g^m=(g^n)^k=e$. Conversely, suppose
$g^m=e$ and write $m=qn+r$ with $0\le r<n$. Then $g^r=g^m\,(g^n)^{-q}=e$, and since $r<n$
the minimality of $n$ forces $r=0$. Hence $n\mid m$.
\end{proof}

\begin{lemma}\label{lem:inv}
Let $G$ be a finite abelian group and $g\in G$. Then $\ord(g^{-1})=\ord(g)$.
\end{lemma}

\begin{proof}
For every $k\ge 0$ we have $(g^{-1})^k=(g^k)^{-1}$, so $(g^{-1})^k=e$ if and only if $g^k=e$.
The least positive $k$ with this property is therefore the same for $g$ and for $g^{-1}$.
\end{proof}

\begin{lemma}\label{lem:selfinv}
Let $g\in G$. Then $g=g^{-1}$ if and only if $g^2=e$.
\end{lemma}

\begin{proof}
Multiply the equation $g=g^{-1}$ on the left by $g$, and conversely multiply $g^2=e$ on the
left by $g^{-1}$.
\end{proof}

\begin{lemma}\label{lem:ord2}
Let $g\in G$. Then $g^2=e$ if and only if $\ord(g)=2$.
\end{lemma}

\begin{proof}
If $\ord(g)=2$ then $g^2=e$ by definition. Conversely, if $g^2=e$ then $\ord(g)$ divides $2$
by Lemma~\ref{lem:div}, so $\ord(g)=2$.
\end{proof}

We now turn to the combinatorial input. An \emph{involution} of a set $X$ is a map
$\sigma\colon X\to X$ with $\sigma\circ\sigma=\mathrm{id}_X$; its set of fixed points is
$\Fix\sigma=\{x\in X:\sigma(x)=x\}$.

\begin{lemma}\label{lem:parity}
Let $X$ be a finite set and $\sigma$ an involution of $X$. Then
$|X|\equiv|\Fix\sigma|\pmod 2$.
\end{lemma}

\begin{proof}
If $x\notin\Fix\sigma$ then also $\sigma(x)\notin\Fix\sigma$, since $\sigma(\sigma(x))=x\neq
\sigma(x)$. Thus $X\setminus\Fix\sigma$ is a union of the sets $\{x,\sigma(x)\}$ with
$x\notin\Fix\sigma$; two such sets are either equal or disjoint, and each has exactly two
elements. If $m$ denotes the number of these sets, then $X\setminus\Fix\sigma$ is the disjoint
union of $m$ two-element sets and $|X|=|\Fix\sigma|+m$. Hence $|X|\equiv|\Fix\sigma|\pmod 2$.
\end{proof}

\begin{corollary}\label{cor:fix}
Let $X$ be a finite set of even cardinality and $\sigma$ an involution of $X$. Then
$|\Fix\sigma|$ is even; in particular, $\sigma$ has at least two fixed points.
\end{corollary}

\begin{proof}
The first assertion is Lemma~\ref{lem:parity}, and an even number of fixed points is at least
two.
\end{proof}

\begin{theorem}\label{thm:main}
Let $G$ be a finite group of even order. Then the number of elements of order $2$ in $G$ is
odd. In particular, $G$ contains an element of order $2$.
\end{theorem}

\begin{proof}
Let $\sigma\colon G\to G$ be the map $\sigma(g)=g^{-1}$. Since $(g^{-1})^{-1}=g$, $\sigma$ is
an involution of $G$, and by Lemma~\ref{lem:selfinv} its fixed point set is
$F=\{g\in G: g^2=e\}$. As $|G|$ is even, $|F|$ is even by Corollary~\ref{cor:fix}. The
identity lies in $F$ and has order $1$, while by Lemma~\ref{lem:ord2} every $g\in F$ with
$g\neq e$ has order $2$; conversely every element of order $2$ lies in $F$. Hence the set of
elements of order $2$ is exactly $F\setminus\{e\}$, which has $|F|-1$ elements, an odd number.
An odd number is nonzero, which gives the last assertion.
\end{proof}

\section{Applications}\label{sec:app}

We use the following two classical results.

\begin{theorem}[Cauchy~\cite{Cauchy}]\label{thm:cauchy}
If a prime $p$ divides $|G|$, then $G$ contains an element of order $p$.
\end{theorem}

\begin{proposition}\label{prop:exp2}
Suppose that $g^2=e$ for every $g\in G$. Then $G$ is abelian and $|G|$ is a power of $2$.
\end{proposition}

\begin{proof}
By Lemma~\ref{lem:selfinv} every element of $G$ equals its own inverse. For $a,b\in G$ we get
$ab=(ab)^{-1}=b^{-1}a^{-1}=ba$, so $G$ is abelian. If an odd prime $p$ divided $|G|$, then by
Theorem~\ref{thm:cauchy} there would be $g\in G$ with $\ord(g)=p$; but $g^2=e$, so $p$ divides
$2$ by Lemma~\ref{lem:div}, which is absurd. Hence $|G|$ has no odd prime divisor, that is,
$|G|=2^k$ for some $k\ge 0$.
\end{proof}

\begin{theorem}[Frobenius~\cite{Frobenius}]\label{thm:frob}
Let $G$ be a finite group and $n$ a positive integer. Then the number of elements $x\in G$
with $x^n=e$ is a multiple of $n$.
\end{theorem}

\begin{corollary}\label{cor:modp}
Let $p$ be a prime dividing $|G|$. Then the number of elements of order $p$ in $G$ is
congruent to $-1$ modulo $p$.
\end{corollary}

\begin{proof}
Let $S=\{x\in G: x^p=e\}$; by Theorem~\ref{thm:frob} with $n=p$, $|S|\equiv 0\pmod p$. By
Lemma~\ref{lem:div}, an element $x\in S$ has $\ord(x)\in\{1,p\}$, and $\ord(x)=1$ exactly when
$x=e$. Conversely every element of order $p$ lies in $S$. So the number of elements of order
$p$ is $|S|-1\equiv-1\pmod p$.
\end{proof}

\begin{remark}
For $p=2$, Corollary~\ref{cor:modp} recovers Theorem~\ref{thm:main}; the point of
Section~2 is that the case $p=2$ needs nothing beyond the pairing argument.
\end{remark}

\begin{thebibliography}{9}
\bibitem{Cauchy} A.-L. Cauchy, \emph{M\'emoire sur les arrangements que l'on peut former
avec des lettres donn\'ees}, Exercices d'analyse et de physique math\'ematique \textbf{3}
(1844), 151--252.
\bibitem{Frobenius} G. Frobenius, \emph{Verallgemeinerung des Sylowschen Satzes},
Sitzungsber. K\"onigl. Preuss. Akad. Wiss. Berlin (1895), 981--993.
\end{thebibliography}

\end{document}
